% संपूर्ण मराठी ग्रंथातील प्रकरण 75: चरित्रे
% हा संपादनयोग्य स्रोतखंड आहे. संकलनासाठी संपूर्ण स्रोतसंग्रहातील
% अक्षररूपे, पूर्वभाग, चिन्हव्याख्या आणि आकृती-स्रोत आवश्यक आहेत.
% मूळ: Open Logic Project; मजकूर आणि रूपांतर CC BY 4.0.
% यंत्रानुवाद, दुरुस्ती आणि तपासणी: OpenAI Codex —
% GPT-5.6 Sol आणि GPT-6 Sol, दोन्ही Ultra effort.
% दुरुस्ती-पुनरावलोकन: GPT-6.1 Sol, Ultra effort.
\chapter{चरित्रे}\label{his:bio:chap}










\section{गेऑर्ग कँटर}\label{his:bio:can:sec}



गेऑर्ग कँटर (\textsc{गे}-ऑर्ग \textsc{कान}-टोर) यांच्या एका जुन्या चरित्रात असा दावा केला होता की रशियातील सेंट पीटर्सबर्गकडे जाणाऱ्या जहाजावर त्यांचा जन्म झाला आणि ते तिथे सापडले; त्यांचे पालक कोण होते हे माहीत नव्हते. मात्र हा दावा खरा नाही. कँटर यांचा जन्म १८४५ मध्ये सेंट पीटर्सबर्ग येथे झाला होता.

कँटर यांनी १८६७ मध्ये बर्लिन विद्यापीठातून गणितातील डॉक्टरेट मिळवली. संच उपपत्तीतील त्यांच्या कार्यासाठी ते ओळखले जातात आणि संशोधनाची स्वतंत्र शाखा म्हणून संच उपपत्तीची स्थापना करण्याचे श्रेय त्यांना दिले जाते. निरनिराळ्या आकारमानांचे अनंत संच असतात हे सर्वप्रथम त्यांनी सिद्ध केले. त्यांच्या उपपत्तींबद्दल, विशेषतः अनंतांच्या उपपत्तीबद्दल, तत्कालीन गणितज्ञांत बरीच चर्चा झाली आणि त्यांचे कार्य वादग्रस्त ठरले.

कँटर यांच्या धार्मिक श्रद्धा आणि गणितीय कार्य यांची घट्ट सांगड होती. अनंतपार संख्यांची उपपत्ती आपल्याला थेट ईश्वराकडून कळली, असा दावाही त्यांनी केला होता. आयुष्याच्या पुढील काळात त्यांना मानसिक आजाराचा त्रास झाला. १८८४ पासून त्यांना रुग्णालयात दाखल करावे लागले; नंतरच्या काळात असे अधिक वारंवार घडले. या स्रोताच्या मांडणीनुसार, त्यांच्या कार्यावरील तीव्र टीका आणि गणितज्ञ लिओपोल्ड क्रोनेकर यांच्याशी झालेला मतभेद यांमुळे त्यांना नैराश्य आले आणि गणितातील रस कमी झाला. नैराश्याच्या काळात कँटर तत्त्वज्ञान आणि साहित्याकडे वळत. शेक्सपियरची नाटके फ्रान्सिस बेकन यांनी लिहिली होती, असा सिद्धांतही त्यांनी प्रकाशित केला.

६ जानेवारी १९१८ रोजी हाल येथील आरोग्यधामात कँटर यांचे निधन झाले.

\paragraph{अधिक वाचन.}
कँटर यांच्या सविस्तर चरित्रांसाठी Dauben (1990) आणि Grattan-Guinness (1971) पाहा. त्यांच्या अपरंपरागत मतांचे वर्णन बीबीसी रेडिओ ४ वरील \emph{A Brief History of Mathematics} (Sautoy 2014) या कार्यक्रमातही आहे. कँटर यांच्या उपपत्ती रॅपच्या रूपात ऐकायच्या असतील, तर Rose (2012) पाहा.













\section{ॲलोन्झो चर्च}\label{his:bio:chu:sec}



ॲलोन्झो चर्च यांचा जन्म १४ जून १९०३ रोजी वॉशिंग्टन, डीसी येथे झाला. लहानपणी हवेच्या बंदुकीशी संबंधित एका अपघातामुळे त्यांचा एक डोळा आंधळा झाला. कनेक्टिकट येथील महाविद्यालयपूर्व शाळेचे शिक्षण त्यांनी १९२० मध्ये पूर्ण केले आणि त्याच वर्षी प्रिन्स्टन विद्यापीठात शिक्षण सुरू केले. १९२७ मध्ये त्यांनी डॉक्टरेट पूर्ण केली. दोन वर्षे परदेशात राहिल्यानंतर चर्च प्रिन्स्टनला परतले. ते अत्यंत नम्र आणि काळजीपूर्वक काम करणारे म्हणून ओळखले जात. फळ्यावरचे त्यांचे लेखन निर्दोष असे. महत्त्वाचे कागद जतन करण्यासाठी ते त्यांवर ड्यूको सिमेंट या पारदर्शक गोंदाचा पातळ थर लावत. शैक्षणिक कामाव्यतिरिक्त त्यांना विज्ञानकथा मासिके वाचायला आवडत; कथांमध्ये काही चूक आढळली तर संपादकांना पत्र लिहायलाही ते कचरत नसत.

चर्च यांचे शैक्षणिक योगदान मोठे होते. त्यांचे विद्यार्थी स्टीफन क्लिनी आणि बार्कली रॉसर यांच्यासह त्यांनी प्रभावी गणनीयतेची एक उपपत्ती, म्हणजे लॅम्डा कलन, विकसित केले. हे काम ॲलन ट्यूरिंग यांनी ट्यूरिंग यंत्र विकसित केले त्या कामापासून स्वतंत्रपणे झाले. गणनीयतेच्या या दोन व्याख्या सममूल्य आहेत. त्यांतून आज \emph{चर्च–ट्यूरिंग प्रबंध} म्हणून ओळखले जाणारे प्रतिपादन पुढे येते: नैसर्गिक संख्यांवरील एखादे फलन प्रभावीपणे गणनीय असते जर आणि फक्त जर ते ट्यूरिंग यंत्राने (किंवा लॅम्डा कलनाने) गणनीय असेल. चर्च यांनी आज \emph{चर्च यांचे प्रमेय} म्हणून ओळखला जाणारा निष्कर्षही सिद्ध केला: प्रथम-क्रम सूत्रांच्या वैधतेची निर्णय समस्या सोडवता येत नाही.

चर्च यांनी वृद्धापकाळापर्यंत काम चालू ठेवले. १९६७ मध्ये त्यांनी प्रिन्स्टन सोडून यूसीएलए येथे प्रवेश केला. १९९० मध्ये निवृत्त होईपर्यंत ते तेथे प्राध्यापक होते. ११ ऑगस्ट १९९५ रोजी वयाच्या ९२व्या वर्षी त्यांचे निधन झाले.

\paragraph{अधिक वाचन.}
चर्च यांचे संक्षिप्त चरित्र Enderton (2019) येथे पाहा. लॅम्डा कलन, चर्च यांचा प्रबंध आणि निर्णय समस्या यांवरील चर्च यांचे मूळ लेख Church (1936; 1936) येथे आहेत. १९३०च्या दशकातील प्रिन्स्टनच्या गणिती समुदायाविषयी चर्च यांची मुलाखत Aspray (1984) यांनी नोंदवली आहे. चर्च यांनी १९३६ ते १९७९ या काळात \emph{Journal of Symbolic Logic} साठी अनेक पुस्तक परीक्षणे लिहिली. ती सर्व जॉन मॅकफार्लेन यांच्या संकेतस्थळावर संग्रहित आहेत (MacFarlane 2015).













\section{गरहार्ड गेंटझेन}\label{his:bio:gen:sec}



गरहार्ड गेंटझेन हे मुख्यतः संरचनात्मक सिद्धता उपपत्तीचे निर्माते म्हणून, आणि विशेषतः नैसर्गिक निगमन व क्रमवर्ती कलन या निष्पत्ती-पद्धती निर्माण केल्याबद्दल ओळखले जातात. त्यांचा जन्म २४ नोव्हेंबर १९०९ रोजी जर्मनीतील ग्राइफ्सवाल्ड येथे झाला. महाविद्यालयपूर्व शाळेत जाण्यापूर्वी गरहार्ड यांचे तीन वर्षे घरीच शिक्षण झाले. शाळेत गेल्यावर शिक्षणाच्या बाबतीत ते आपल्या बहुतेक सहाध्यायांच्या मागे होते. तरीही ते उत्तम विद्यार्थी होते आणि गणितातील त्यांची विशेष क्षमता दिसून येत होती. त्यांची आवड विविध विषयांत होती. उदाहरणार्थ, त्यांनी आपल्या आईसाठी कविता आणि शाळेच्या नाट्यगृहासाठी नाटकेही लिहिली.

गेंटझेन यांनी विद्यापीठीय शिक्षणाची सुरुवात ग्राइफ्सवाल्ड विद्यापीठात केली; त्यानंतर ते गॉटिंगेन (G\"{o}ttingen), म्यूनिख आणि बर्लिन येथे गेले. १९३३ मध्ये त्यांनी गॉटिंगेन (G\"{o}ttingen) विद्यापीठातून हेरमान वाइल यांच्या मार्गदर्शनाखाली डॉक्टरेट मिळवली. (त्यांच्या बहुतांश कामाचे मार्गदर्शन पॉल बर्नायस यांनी केले होते; मात्र नाझींनी त्यांना विद्यापीठातून काढून टाकले.) १९३४ मध्ये गेंटझेन डेव्हिड हिल्बर्ट यांचे सहाय्यक म्हणून काम करू लागले. क्रमवर्ती कलन आणि नैसर्गिक निगमन या निष्पत्ती-पद्धतींवरील त्यांचे \emph{Untersuchungen \"{u}ber das logische Schlie\ss en I--II
  [Investigations Into Logical Deduction I--II]} हे लेख १९३५ मध्ये प्रकाशित झाले. १९३६ मध्ये त्यांनी पेआनो स्वयंसिद्धकांची सुसंगती सिद्ध केली.

गेंटझेन यांचा नाझींशी असलेला संबंध गुंतागुंतीचा होता. त्यांचे मार्गदर्शक बर्नायस यांना जर्मनी सोडण्यास भाग पाडले जात असतानाच गेंटझेन नाझींची निमलष्करी संघटना असलेल्या एसएच्या विद्यापीठीय शाखेत सामील झाले. अनेक जर्मनांप्रमाणे ते नाझी पक्षाचेही सदस्य होते. युद्धकाळात त्यांनी हवाई गुप्तचर विभागात दूरसंचार अधिकारी म्हणून काम केले. मात्र १९४२ मध्ये मानसिक आरोग्य अचानक खालावल्याने त्यांना सेवेतून मुक्त करण्यात आले. गेंटझेन यांची निष्ठा नाझी पक्षाशी होती की शैक्षणिक प्रगती सुनिश्चित करण्यासाठी ते पक्षात सामील झाले, हे स्पष्ट नाही.

१९४३ मध्ये गेंटझेन यांना प्राग येथील जर्मन विद्यापीठाच्या गणित संस्थेत शैक्षणिक पद देऊ करण्यात आले; त्यांनी ते स्वीकारले. मात्र १९४५ मध्ये प्रागच्या नागरिकांनी जर्मन अधिपत्याविरुद्ध बंड केले. सोव्हिएत सैन्य शहरात दाखल झाले आणि विद्यापीठातील सर्व प्राध्यापकांना अटक केली, असे या स्रोतामध्ये म्हटले आहे. नाझी संघटनांतील सदस्यत्वामुळे गेंटझेन यांना सक्तीच्या श्रमांच्या छावणीत नेण्यात आले. ४ ऑगस्ट १९४५ रोजी वयाच्या ३५व्या वर्षी कारागृहातील त्यांच्या कोठडीत कुपोषणामुळे त्यांचे निधन झाले.

\paragraph{अधिक वाचन.}
गेंटझेन यांचे संपूर्ण चरित्र Menzler-Trott (2007) येथे पाहा. नाझी राजवटीखालील गणितज्ञांविषयीचे Segal (2014) यांचे पुस्तकही वाचनीय आहे; त्यात गेंटझेन यांच्या आयुष्याविषयी एक संक्षिप्त नोंद आहे. तार्किक निगमनावरील गेंटझेन यांचे मूळ जर्मन लेख (Gentzen 1935; Gentzen 1935) येथे उपलब्ध आहेत. त्यांच्या लेखांची इंग्रजी भाषांतरे Szabo (1969) यांनी एका खंडात एकत्र केली आहेत; त्यात त्यांचे संक्षिप्त चरित्रही आहे.














\section{कुर्ट गोडेल}\label{his:bio:god:sec}



कुर्ट गोडेल (\textsc{गर}-डल) यांचा जन्म २८ एप्रिल १९०६ रोजी ऑस्ट्रो-हंगेरी साम्राज्यातील ब्र्युन येथे झाला; ते ठिकाण आता झेक प्रजासत्ताकातील ब्र्नो आहे. लहानपणी ते जिज्ञासू आणि हुशार असल्याने घरचे त्यांना “Der kleine Herr Warum” (लहानसे 'का-महोदय') असे म्हणत. प्राथमिक शाळेपासूनच त्यांची शैक्षणिक कामगिरी उत्कृष्ट होती; सर्वोच्च गुणांपेक्षा कमी गुण त्यांना फक्त गणितात मिळाले. तब्येत बरी नसल्याने गोडेल अनेकदा शाळेत गैरहजर राहत आणि त्यांना शारीरिक शिक्षणातून सूट मिळाली होती. बालपणी त्यांना संधिवातजन्य तापाचे निदान झाले. वैद्यकीय तपासणीत त्याच्या उलट निष्कर्ष आला असतानाही, त्या आजाराचा आपल्या हृदयावर कायमचा परिणाम झाला आहे, असे त्यांना आयुष्यभर वाटत राहिले.

गोडेल यांनी १९२४ मध्ये व्हिएन्ना विद्यापीठात शिक्षण सुरू केले आणि १९२९ मध्ये डॉक्टरेटचा प्रबंध पूर्ण केला; डॉक्टरेटची पदवी १९३० मध्ये मिळाली. सुरुवातीला त्यांचा भौतिकशास्त्र शिकण्याचा विचार होता; पण काही अंशी तत्त्वज्ञ रुडॉल्फ कार्नाप यांच्या प्रभावामुळे त्यांची आवड गणिताकडे, विशेषतः तर्कशास्त्राकडे वळली. हान्स हान यांच्या मार्गदर्शनाखाली लिहिलेल्या प्रबंधात त्यांनी समानतेसह प्रथम-क्रम विधेय तर्कशास्त्राचे संपूर्णता प्रमेय सिद्ध केले (G\"odel 1929). या प्रबंधानंतर अवघ्या वर्षभरात त्यांनी आपले सर्वात प्रसिद्ध निष्कर्ष, पहिले आणि दुसरे अपूर्णता प्रमेय, मिळवले (हे G{\"{o}}del 1931 मध्ये प्रसिद्ध झाले). व्हिएन्नातील वास्तव्यात गोडेल 'व्हिएन्ना मंडळ' या विज्ञानाभिमुख तत्त्वज्ञांच्या समूहात सक्रिय होते. या समूहातील कार्नाप यांच्या कार्यावर गोडेल यांच्या निष्कर्षांचा विशेष प्रभाव पडला.

१९३८ मध्ये गोडेल यांनी अडेल निंबुर्स्की यांच्याशी विवाह केला. त्यांचे पालक या विवाहावर नाराज होते: अडेल त्यांच्यापेक्षा सहा वर्षांनी मोठ्या होत्या, त्यांचा आधी घटस्फोट झाला होता आणि त्या एका नाइटक्लबमध्ये नर्तकी म्हणून काम करत होत्या. मात्र सामाजिक दबावाचा गोडेल यांच्यावर परिणाम झाला नाही; त्यांच्या निधनापर्यंत त्यांचा विवाह आनंदी राहिला, असे स्रोत सांगतो.

१९३८ मध्ये नाझी जर्मनीने ऑस्ट्रियाचा विलय केल्यानंतर गोडेल आणि अडेल अमेरिकेत स्थलांतरित झाले. न्यू जर्सीतील प्रिन्स्टन येथील इन्स्टिट्यूट फॉर ॲडव्हान्स्ड स्टडी येथे गोडेल यांनी पद स्वीकारले. ते अंतर्मुख आणि काहीसे विलक्षण स्वभावाचे असले, तरी प्रिन्स्टनमधील त्यांचा काळ सहकार्यपूर्ण आणि फलदायी ठरला. त्यांनी संच उपपत्ती, तत्त्वज्ञान आणि भौतिकशास्त्र यांवर निबंध प्रकाशित केले. याच संस्थेतील सहकारी अल्बर्ट आइन्स्टाइन यांच्याशी त्यांची विशेष घनिष्ठ मैत्री झाली.

आयुष्याच्या उत्तरार्धात गोडेल यांचे मानसिक आरोग्य खालावले. १९७७ मध्ये त्यांच्या पत्नीला रुग्णालयात दाखल करावे लागल्याने त्या त्यांच्यासाठी जेवण बनवू शकत नव्हत्या. गोडेल यांना आयुष्यभर मानसिक आरोग्याच्या समस्या होत्या; त्यातून त्यांच्या मनात संशय बळावला. आपल्याला विष घातले जाईल या तीव्र भीतीने त्यांनी अन्न घेणे बंद केले. १४ जानेवारी १९७८ रोजी प्रिन्स्टन येथे उपासमारीने त्यांचे निधन झाले, असे या स्रोताचे वर्णन आहे.

\paragraph{अधिक वाचन.}
गोडेल यांचे सविस्तर चरित्र John Dawson (1997) येथे पाहा. त्यांच्या चरित्राविषयीचे आणखी लेख आणि तर्कशास्त्र व तत्त्वज्ञानातील योगदानावरील निबंध Wang (1990), Baaz (2011), Takeuti (2003) आणि Sigmund (2007) येथे मिळतील.

गोडेल यांचा डॉक्टरेट प्रबंध मूळ जर्मन भाषेत (G\"odel 1929) येथे उपलब्ध आहे. अपूर्णता प्रमेयांचा मूळ लेख (G{\"{o}}del 1931) येथे आहे. त्यांचे प्रकाशित व अप्रकाशित सर्व लेखन आणि निवडक पत्रव्यवहार यांची इंग्रजी आवृत्ती त्यांच्या \emph{Collected Papers} Feferman (1986; 1990) मध्ये उपलब्ध आहे.

गोडेल यांच्या अपूर्णता प्रमेयांच्या सविस्तर विवेचनासाठी Smith (2013) पाहा. या प्रमेयांवरील अनौपचारिक तात्त्विक चर्चेसाठी मार्क लिन्झेनमायर यांचे पॉडकास्ट (Linsenmayer 2014) पाहा.













\section{एमी नोएथर}\label{his:bio:noe:sec}



एमी नोएथर (\textsc{नर}-टर) यांचा जन्म २३ मार्च १८८२ रोजी जर्मनीतील एरलांगेन येथे एका उच्च मध्यमवर्गीय विद्वान कुटुंबात झाला. त्यांना “आधुनिक बीजगणिताची जननी” म्हटले जाते. स्त्रियांच्या शिक्षणातील मोठ्या अडथळ्यांवर मात करत त्यांनी गणित आणि भौतिकशास्त्र या दोन्ही क्षेत्रांत मूलगामी योगदान दिले. त्या काळी जर्मनीतील मुलींनी कलांचे शिक्षण घ्यावे अशी अपेक्षा होती आणि त्यांना महाविद्यालयपूर्व शाळांमध्ये प्रवेश नव्हता. तरीही नोएथर यांनी गॉटिंगेन आणि एरलांगेन विद्यापीठांतील वर्गांत श्रोती म्हणून उपस्थिती लावली; एरलांगेन येथे त्यांचे वडील गणिताचे प्राध्यापक होते. अखेरीस १९०४ मध्ये महिला विद्यार्थिनींना प्रवेश देण्याचे एरलांगेनचे धोरण बदलल्यावर त्या तेथे विद्यार्थिनी म्हणून दाखल झाल्या. १९०७ मध्ये त्यांनी गणितात डॉक्टरेट मिळवली.

पात्रता असूनही नोएथर यांना कारकिर्दीत मोठा विरोध सहन करावा लागला. १९०८ ते १९१५ या काळात त्यांनी एरलांगेन येथे विनावेतन अध्यापन केले. याच काळात तत्कालीन अग्रगण्य गणितज्ञ डेव्हिड हिल्बर्ट यांचे त्यांच्याकडे लक्ष गेले आणि त्यांनी नोएथर यांना गॉटिंगेनला बोलावले. मात्र स्त्रियांना प्राध्यापकपदे मिळण्यास बंदी होती. त्यामुळे त्यांना हिल्बर्ट यांच्या नावाखालीच व्याख्याने देता आली, तीही विनावेतन. याच काळात त्यांनी आज 'नोएथर यांचे प्रमेय' म्हणून ओळखला जाणारा निष्कर्ष सिद्ध केला; सैद्धांतिक भौतिकशास्त्रात आजही त्याचा उपयोग होतो. अखेरीस १९१९ मध्ये त्यांना अध्यापनाचा अधिकार मिळाला. विद्यापीठातील सहकाऱ्यांचा विरोध सुरूच असताना हिल्बर्ट यांनी “महाशयांनो, विद्याशाखेची सभा स्नानगृह नव्हे,” असे उत्तर दिल्याचे सांगितले जाते.

१९२०च्या दशकाच्या उत्तरार्धात नोएथर यांनी अमूर्त बीजगणितावर लक्ष केंद्रित केले आणि त्यांच्या योगदानाने या क्षेत्रात आमूलाग्र बदल घडवला. आपल्या सिद्धतांमध्ये त्या अनेकदा तथाकथित आरोही साखळी अट वापरत. या अटीनुसार विशिष्ट संचांची काटेकोरपणे वाढत जाणारी अनंत साखळी नसते. उदाहरणार्थ, आता नोएथरियन वलये म्हणून ओळखल्या जाणाऱ्या काही बीजगणितीय रचनांत आदर्शांचा $I_1
\subsetneq I_2 \subsetneq \dots$ असा अनंत अनुक्रम नसतो. ही अट कोणत्याही अंशतः क्रमावर व्यापक करता येते (बीजगणितात आदर्शांना उपसंच-संबंधाने लावलेला क्रम हे तिचे विशेष प्रकरण आहे). तिची द्वैती अवरोही साखळी अटही विचारात घेता येते: अंशतः क्रमातील प्रत्येक काटेकोरपणे \emph{उतरता} अनुक्रम अखेरीस थांबतो. एखादा अंशतः क्रम ही अवरोही साखळी अट पूर्ण करत असेल, तर $<$ या $\Nat$ वरील क्रमाने जसे विगमन करता येते, तसे त्या क्रमानेही विगमन करता येते. अशा क्रमांना \emph{सुस्थापित} किंवा \emph{नोएथरियन} म्हणतात; त्यानुसार मिळणाऱ्या सिद्धता-तत्त्वाला \emph{नोएथरियन विगमन} म्हणतात.

नोएथर ज्यू होत्या. १९३३ मध्ये नाझी सत्तेवर आल्यावर त्यांना पदावरून काढून टाकण्यात आले. सुदैवाने त्यांना अमेरिकेतील पेनसिल्व्हेनिया राज्यातील ब्रिन मॉर येथे तात्पुरत्या पदासाठी स्थलांतर करता आले. त्या काळात त्यांनी प्रिन्स्टनमध्येही व्याख्याने दिली; मात्र ते विद्यापीठ स्त्रियांना अनुकूल नाही, असे त्यांना वाटले (Dick 1981, पृ.~81). १९३५ मध्ये त्यांच्या गर्भाशयातील गाठ काढण्यासाठी शस्त्रक्रिया झाली. शस्त्रक्रियेनंतर झालेल्या संसर्गामुळे त्यांचे निधन झाले आणि ब्रिन मॉर येथे त्यांचे दफन झाले.

\paragraph{अधिक वाचन.}
नोएथर यांच्या चरित्रासाठी Dick (1981) पाहा. पेरिमीटर इन्स्टिट्यूट फॉर थिअरेटिकल फिजिक्सची त्यांच्या जीवनावरील आणि प्रभावावरील व्याख्याने ऑनलाइन उपलब्ध आहेत (Institute 2015). वाचनाऐवजी ऐकायचे असल्यास, \emph{Stuff You Missed in History Class} या मालिकेतील नोएथर यांच्या जीवनावरील आणि प्रभावावरील पॉडकास्ट (Frey 2015) पाहा. नोएथर यांचे संकलित लेखन मूळ जर्मन भाषेत (Jacobson 1983) येथे उपलब्ध आहे.













\section{रोझा पेटर}\label{his:bio:pet:sec}



रोझा पेटर यांचा जन्म १७ फेब्रुवारी १९०५ रोजी हंगेरीतील बुडापेस्ट येथे रोझा पोलित्झर या नावाने झाला. पुनरावर्ती फलनांवरील कार्यासाठी त्या विशेष ओळखल्या जातात; पुनरावर्तन उपपत्ती हे क्षेत्र निर्माण होण्यात हे कार्य अत्यावश्यक ठरले.

पेटर यांचे बालपण राजकीय उलथापालथीच्या काळात गेले; त्या किशोरवयीन असताना पहिले महायुद्ध सुरू होते. तरीही त्यांना बुडापेस्टमधील प्रतिष्ठित मारिया तेरेझिया मुलींच्या शाळेत शिकता आले आणि १९२२ मध्ये त्यांनी तेथील शिक्षण पूर्ण केले. त्यानंतर त्यांनी बुडापेस्टच्या पाझमान्य पेटर विद्यापीठात शिक्षण घेतले; पुढे त्याचे नाव लोरांड एटव्होश विद्यापीठ झाले. वडिलांच्या आग्रहामुळे त्यांनी सुरुवातीला रसायनशास्त्र घेतले, पण नंतर गणिताकडे वळल्या आणि १९२७ मध्ये पदवीधर झाल्या. माध्यमिक शाळेत गणित शिकवण्याची पात्रता असूनही त्या काळची आर्थिक परिस्थिती भीषण होती; महामंदीने जागतिक अर्थव्यवस्था प्रभावित झाली होती. म्हणून त्यांनी शिकवण्या आणि खाजगी गणित अध्यापन अशी प्रसंगोपात्त कामे केली. अखेरीस गणितातील पदव्युत्तर शिक्षणासाठी त्या पुन्हा विद्यापीठात आल्या. सुरुवातीला त्यांचा संख्या उपपत्तीवर काम करण्याचा विचार होता; पण आपले निष्कर्ष आधीच सिद्ध झाले आहेत हे कळल्यावर त्यांनी गणितच सोडून देण्याचा जवळजवळ निर्णय घेतला. त्यांना गोडेल यांच्या अपूर्णता प्रमेयांवर काम करण्यास प्रोत्साहन मिळाले आणि नकळत त्यांनी त्यांचे काही निष्कर्ष वेगळ्या पद्धतीने सिद्ध केले. त्यामुळे त्यांचा आत्मविश्वास परत आला. डेव्हिड हिल्बर्ट यांच्या पायाभूत कार्यक्रमातून प्रेरणा घेऊन पेटर यांनी पुनरावर्तन उपपत्तीवरील आपले पहिले लेख लिहिले. १९३५ मध्ये त्यांनी डॉक्टरेट मिळवली आणि १९३७ मध्ये त्या \emph{Journal of
  Symbolic Logic} च्या संपादक झाल्या.

पेटर यांचे आरंभीचे लेख पुनरावर्ती फलनांच्या उपपत्तीचे संस्थापक योगदान म्हणून व्यापकपणे मानले जातात. (P{\'e}ter 1935a) मध्ये त्यांनी पुनरावर्तनाच्या विविध प्रकारांमधील संबंध तपासले. (P\'eter 1935b) मध्ये त्यांनी पुनरावर्ती पद्धतीने परिभाषित केलेले एक विशिष्ट फलन आदिम पुनरावर्ती नाही, हे दाखवले. यामुळे विल्हेल्म ॲकरमन यांचा आधीचा निष्कर्ष सोपा झाला. पेटर यांचे हे सुलभ केलेले फलन आज अनेकदा ॲकरमन फलन म्हणतात; काही वेळा अधिक अचूकपणे ॲकरमन–पेटर फलनही म्हणतात. पुनरावर्ती फलनांच्या उपपत्तीवरील पहिले पुस्तक त्यांनी लिहिले (P\'eter 1951).

त्यांच्या कार्याचे महत्त्व आणि प्रभाव मोठा असूनही पेटर यांना १९४५ पर्यंत पूर्णवेळ अध्यापनाचे पद मिळाले नाही. दुसऱ्या महायुद्धात नाझींनी हंगेरीवर कब्जा केला असताना यहुदीविरोधी कायद्यांमुळे त्यांना शिकवण्याची परवानगी नव्हती. १९४४ मध्ये सरकारने बुडापेस्टमध्ये ज्यू लोकांसाठी बंदिस्त वस्ती निर्माण केली; ती शहराच्या इतर भागांपासून अलग ठेवली गेली होती आणि सशस्त्र रक्षक तिच्यावर पहारा देत होते. १९४५ मध्ये ही वस्ती मुक्त होईपर्यंत पेटर यांना तेथे राहण्यास भाग पाडले गेले. त्यानंतर त्यांनी बुडापेस्ट टीचर्स ट्रेनिंग कॉलेजमध्ये आणि १९५५ पासून एटव्होश लोरांड विद्यापीठात अध्यापन केले. ‘गणितातील अकादमिक डॉक्टर’ हा दर्जा मिळवणाऱ्या त्या पहिल्या हंगेरियन महिला गणितज्ञ आणि हंगेरियन विज्ञान अकादमीवर निवडून जाणाऱ्या पहिल्या महिला होत्या.

पेटर या गणित उत्कटतेने शिकवणाऱ्या अध्यापिका म्हणून ओळखल्या जात. केवळ व्याख्यान देण्याऐवजी विद्यार्थ्यांसोबत गणिती समस्यांचे स्वरूप आणि सौंदर्य शोधणे त्यांना अधिक आवडत असे. म्हणून विद्यार्थी त्यांना प्रेमाने “रोझा मावशी” म्हणत. १९७७ मध्ये वयाच्या ७१व्या वर्षी त्यांचे निधन झाले.

\paragraph{अधिक वाचन.}
अधिक चरित्रात्मक वाचनासाठी (O'Connor 2014) आणि (Andr\'asfai 1986) पाहा. पेटर यांची एक छोटी मुलाखत Tamassy (1994) यांनी घेतली होती. गणिताविषयी रंजक वाचनासाठी पेटर यांचे \emph{Playing With Infinity} हे पुस्तक (P\'eter 2010) पाहा.












\section{ज्युलिया रॉबिन्सन}\label{his:bio:rob:sec}



ज्युलिया बोमन रॉबिन्सन या अमेरिकन गणितज्ञ होत्या. त्या मुख्यतः निर्णय समस्यांवरील कार्यासाठी, आणि विशेषतः हिल्बर्ट यांच्या दहाव्या समस्येच्या समाधानातील योगदानासाठी ओळखल्या जातात. रॉबिन्सन यांचा जन्म ८ डिसेंबर १९१९ रोजी मिसुरी राज्यातील सेंट लुईस येथे झाला. लहानपणापासूनच संख्यांविषयी कुतूहल होते, अशी त्यांची आठवण आहे (Reid 1986, पृ.~4). वयाच्या नवव्या वर्षी त्यांना स्कार्लेट फीव्हर झाला आणि नंतर संधिवातजन्य तापाचे अनेक झटके आले. त्यामुळे त्यांना बराच वेळ अंथरुणात घालवावा लागला आणि शिक्षणात त्या मागे पडल्या. खाजगी शिक्षकांच्या मदतीने त्यांनी ही तूट भरून काढली; मात्र आजाराचे शारीरिक परिणाम त्यांच्या पुढील आयुष्यावरही राहिले.

बालपणीच्या अडचणी असूनही रॉबिन्सन यांनी गणित आणि विज्ञानातील अनेक पुरस्कारांसह माध्यमिक शिक्षण पूर्ण केले. विद्यापीठीय शिक्षणाची सुरुवात त्यांनी सॅन डिएगो स्टेट कॉलेजमध्ये केली आणि शेवटच्या वर्षी कॅलिफोर्निया विद्यापीठ, बर्कली येथे प्रवेश घेतला. तेथे गणितज्ञ राफाएल रॉबिन्सन यांचा त्यांच्यावर प्रभाव पडला. दोघांची मैत्री झाली आणि १९४१ मध्ये त्यांनी विवाह केला. विद्याशाखेतील सदस्याची पत्नी असल्याने ज्युलिया यांना बर्कलीच्या गणित विभागात अध्यापनाची परवानगी नव्हती. त्या गणिताचे वर्ग श्रोती म्हणून घेत राहिल्या, पण विद्यापीठ सोडून कुटुंब सुरू करण्याची त्यांची इच्छा होती. विवाहानंतर थोड्याच काळात त्यांना न्यूमोनिया झाला. बालपणीच्या संधिवातजन्य तापामुळे त्यांच्या हृदयावर बरीच व्रणयुक्त ऊतकांची वाढ झाली आहे, असे त्यांना सांगण्यात आले. या स्थितीची तीव्रता पाहून डॉक्टरांनी त्या चाळीस वर्षांपलीकडे जगणार नाहीत असा अंदाज व्यक्त केला आणि त्यांना मुले न होऊ देण्याचा सल्ला दिला (Reid 1986, पृ.~13).

रॉबिन्सन दीर्घकाळ नैराश्यात होत्या; अखेरीस त्यांनी गणिताचे शिक्षण सुरू ठेवायचे ठरवले. त्या बर्कलीला परतल्या आणि ॲल्फ्रेड टार्स्की यांच्या मार्गदर्शनाखाली १९४८ मध्ये डॉक्टरेट पूर्ण केली. वास्तविक संख्यांची प्रथम-क्रम उपपत्ती निर्णेय आहे, हे टार्स्की यांनी दाखवले होते; गोडेल यांच्या कार्यावरून नैसर्गिक संख्यांची प्रथम-क्रम उपपत्ती अनिर्णेय आहे, हे निष्पन्न होत होते. परिमेय संख्यांची प्रथम-क्रम उपपत्ती निर्णेय आहे की नाही, हा मोठा अनिर्णित प्रश्न होता. आपल्या प्रबंधात (1949) रॉबिन्सन यांनी ती अनिर्णेय आहे, हे सिद्ध केले.

निर्णय समस्यांत रस असल्याने रॉबिन्सन यांनी पुढे हिल्बर्ट यांच्या दहाव्या समस्येचे समाधान शोधण्याचा प्रयत्न केला. १९०० मध्ये डेव्हिड हिल्बर्ट यांनी मांडलेल्या २३ प्रसिद्ध गणिती समस्यांपैकी ही एक होती. पूर्णांक गुणांक असलेल्या $3x^2 - 2y +3 = 0$ यांसारख्या बहुपदी समीकरणाला पूर्णांकांमध्ये उत्तर आहे की नाही, हे मर्यादित वेळेत सांगणारा कलनविधी आहे का, असा दहाव्या समस्येचा प्रश्न आहे. अशा प्रश्नांना \emph{डायोफँटाइन समस्या} म्हणतात. काही सुरुवातीच्या यशांनंतर रॉबिन्सन यांनी याच समस्येवर काम करणारे मार्टिन डेव्हिस आणि हिलरी पुटनम यांच्यासोबत काम सुरू केले. अज्ञात राशी घातांकांतही येऊ शकतात अशा घातांकीय डायोफँटाइन समस्या अनिर्णेय आहेत, हे त्यांनी दाखवले. तसेच पुढे “J.R.” म्हणून ओळखल्या गेलेल्या एका विशिष्ट गृहीतकावरून हिल्बर्ट यांची दहावी समस्या अनिर्णेय आहे हे निष्पन्न होते, असेही त्यांनी सिद्ध केले (Davis 1961). १९६०च्या दशकातही रॉबिन्सन या समस्येवर काम करत राहिल्या. अखेरीस १९७० मध्ये तरुण रशियन गणितज्ञ युरी मातियासेविच यांनी J.R. गृहीतक सिद्ध केले. एकत्रित निष्कर्षाला आता मातियासेविच–रॉबिन्सन–डेव्हिस–पुटनम प्रमेय, संक्षेपात MRDP प्रमेय, म्हणतात. मातियासेविच आणि रॉबिन्सन यांची मैत्री झाली आणि त्यांनी अनेक लेखांवर एकत्र काम केले. मातियासेविच यांना लिहिलेल्या एका पत्रात रॉबिन्सन यांनी असे म्हटले होते: “प्रत्यक्षात, हजारो मैल दूर राहूनही एकत्र काम केल्यामुळे आपल्यापैकी कोणीही एकट्याने जितकी प्रगती केली असती त्याहून अधिक प्रगती आपण करत आहोत, याचा मला खूप आनंद आहे” (Matijasevich 1992, पृ.~45).

रॉबिन्सन या अमेरिकन मॅथेमॅटिकल सोसायटीच्या पहिल्या महिला अध्यक्ष होत्या. राष्ट्रीय विज्ञान अकादमीच्या गणित विभागात निवडून जाणाऱ्या त्या पहिल्या महिला होत्या. रक्ताच्या कर्करोगाचे निदान झाल्यानंतर ३० जुलै १९८५ रोजी वयाच्या ६५व्या वर्षी त्यांचे निधन झाले.

\paragraph{अधिक वाचन.}
रॉबिन्सन यांचे गणितीय लेख त्यांच्या \textit{Collected Works} (Robinson 1996) मध्ये उपलब्ध आहेत. त्यात राष्ट्रीय विज्ञान अकादमीसाठी लिहिलेल्या त्यांच्या चरित्रात्मक स्मृतिलेखाचे पुनर्मुद्रणही आहे (Feferman 1994). रॉबिन्सन यांची मोठी बहीण कॉन्स्टन्स रीड यांनी मुलाखतींवर आधारित “Autobiography of Julia” (Reid 1986) आणि त्यांचे सविस्तर चरित्र (Reid 1996) प्रकाशित केले. रॉबिन्सन आणि हिल्बर्ट यांच्या दहाव्या समस्येवरील लघुपटाचे दिग्दर्शन जॉर्ज चिचेरी यांनी केले (Csicsery 2016). मातियासेविच यांचे रॉबिन्सन यांच्यासोबतचे सहकार्य आणि त्यांच्या कार्यावरील रॉबिन्सन यांचा प्रभाव यांविषयीच्या संक्षिप्त स्मृतिलेखासाठी (Matijasevich 1992) पाहा.














\section{बर्ट्रंड रसेल}\label{his:bio:rus:sec}



बर्ट्रंड रसेल यांना आधुनिक विश्लेषणात्मक तत्त्वज्ञानाच्या संस्थापकांपैकी एक मानले जाते. १८ मे १८७२ रोजी जन्मलेले रसेल तत्त्वज्ञान आणि तर्कशास्त्रातील कामासाठी तर ओळखले जातच; याशिवाय त्यांनी अनेक विषयांवर सर्वसामान्य वाचकांसाठी पुस्तके लिहिली. आयुष्यभर ते सक्रिय राजकीय कार्यकर्तेही होते.

रसेल यांचा जन्म वेल्समधील मॉनमथशरच्या ट्रेलेच येथे झाला. त्यांचे आईवडील ब्रिटिश उमराव वर्गातील होते. ते स्वतंत्र विचारांचे होते आणि त्या काळी बोस्टनमधील क्रांतिकारी विचारांच्या लोकांशीही त्यांची मैत्री होती. दुर्दैवाने रसेल लहान असतानाच त्यांच्या आईवडिलांचे निधन झाले; म्हणून ते आजीआजोबांकडे राहायला गेले. तेथे त्यांचे धार्मिक संस्कारांत संगोपन झाले, जे त्यांच्या आईवडिलांना कोणत्याही परिस्थितीत टाळायचे होते. नैतिकतेच्या सर्व बाबतीत त्यांची आजी अतिशय कडक होती. किशोरवयात त्यांचे बहुतांश शिक्षण खाजगी शिक्षकांकडून घरीच झाले.

विश्लेषणात्मक तत्त्वज्ञानावर, विशेषतः तर्कशास्त्रावर, रसेल यांचा प्रभाव फार मोठा आहे. त्यांनी केंब्रिजच्या ट्रिनिटी कॉलेजमध्ये गणित आणि तत्त्वज्ञान शिकले; तेथे गणितज्ञ आणि तत्त्वज्ञ ॲल्फ्रेड नॉर्थ व्हाइटहेड यांचा त्यांच्यावर प्रभाव पडला. १९१० मध्ये रसेल आणि व्हाइटहेड यांनी \emph{Principia Mathematica} चा पहिला खंड प्रकाशित केला. त्यात गणिताचे तर्कशास्त्रात न्यूनीकरण करता येते, या भूमिकेचे त्यांनी समर्थन केले. पुढे रसेल यांनी शेकडो पुस्तके, निबंध आणि राजकीय पत्रके प्रकाशित केली. १९५० मध्ये त्यांना साहित्यातील नोबेल पारितोषिक मिळाले.

रसेल राजकारणात आणि सामाजिक चळवळींत खोलवर सहभागी होते. पहिल्या महायुद्धात शांततावादी कृती आणि निषेध यांमुळे त्यांना अटक झाली आणि सहा महिन्यांच्या तुरुंगवासाची शिक्षा झाली. तुरुंगात त्यांना वाचता आणि लिहिता येत असे; हा अनुभव “बराच सुखावह” वाटल्याचा त्यांचा दावा होता. पुढेही त्यांनी शांततेसाठी काम केले, मात्र सर्व परिस्थितींमध्ये युद्धाला विरोध करावा अशी त्यांची भूमिका नव्हती. १९६१ मध्ये अण्वस्त्रनिःशस्त्रीकरणाच्या सभेत सहभागी झाल्यामुळे त्यांना पुन्हा तुरुंगवास झाला. १९४८ मध्ये ते विमान अपघातातूनही बचावले. या स्रोताच्या कथनानुसार, तेव्हा धूम्रपानासाठी राखीव विभागात बसलेले लोकच बचावले; म्हणून धूम्रपानामुळेच आपला जीव वाचला, असा रसेल यांचा दावा होता. रसेल यांनी चार विवाह केले, पण विवाहबाह्य संबंध ठेवत असल्याची त्यांची ख्याती होती. २ फेब्रुवारी १९७० रोजी वयाच्या ९७व्या वर्षी वेल्समधील पेनरिनड्यूड्रेथ येथे त्यांचे निधन झाले.

\paragraph{अधिक वाचन.}
रसेल यांनी १८७२ ते १९६७ या कालखंडातील आपले जीवन सांगणारे तीन भागांचे आत्मचरित्र लिहिले (Russell 1967; Russell 1968; Russell 1969). मॅकमास्टर विद्यापीठातील बर्ट्रंड रसेल संशोधन केंद्रात त्यांच्या संग्रहित दस्तऐवजांचा ठेवा आहे. त्यांच्या संकलित लेखनाचे खंड, शोधता येणाऱ्या अनुक्रमणिका आणि संग्रहण प्रकल्पांविषयी केंद्राच्या संकेतस्थळावर Duncan (2015) माहिती आहे. रसेल यांचा \emph{On
  Denoting} (Russell 1905) हा लेख विसाव्या शतकातील विश्लेषणात्मक तत्त्वज्ञानातील अभिजात लेख मानला जातो.

रसेल यांच्यावरील स्टॅनफर्ड एन्सायक्लोपीडिया ऑफ फिलॉसॉफीमधील नोंदीत (Irvine 2015) इच्छा आणि राजकीय सिद्धांताविषयी रसेल बोलत असल्याच्या ध्वनिमुद्रिका आहेत. रसेल यांच्या अनेक चित्रित मुलाखती ऑनलाइन उपलब्ध आहेत. उदाहरणार्थ, धूम्रपान आणि विमान अपघातातील सहभागाविषयी त्यांना बोलताना पाहण्यासाठी Russell (n.d.) पाहा. त्यांच्या काही पुस्तकांची, त्यांत \emph{Introduction to Mathematical Philosophy} चा समावेश आहे, विनामूल्य ध्वनिपुस्तके LibriVox (n.d.) येथे उपलब्ध आहेत.













\section{ॲल्फ्रेड टार्स्की}\label{his:bio:tar:sec}



ॲल्फ्रेड टार्स्की यांचा जन्म १४ जानेवारी १९०१ रोजी पोलंडमधील वॉर्सा येथे झाला; तेव्हा ते शहर रशियन साम्राज्याचा भाग होते. टार्स्की यांचे वर्णन “नेपोलियनसारखे” असे केले जात असे: ते उत्साही, बोलके आणि तीव्र स्वभावाचे होते. त्यांची ऊर्जा व्याख्यानांतही दिसत असे. एकदा व्याख्यानादरम्यान सिगारेट टाकताना त्यांनी कचरापेटीला आग लावली आणि त्यानंतर त्या इमारतीत त्यांना पुन्हा व्याख्यान देण्यास बंदी घालण्यात आली.

लहानपणापासूनच टार्स्की यांना ज्ञानाची तीव्र ओढ होती. आयुष्याच्या पुढील काळात ते विद्यार्थ्यांना सांगत की तर्कशास्त्रात त्यांना 'B' गुण मिळाले म्हणून त्यांनी तो विषय शिकला; तोच त्यांचा 'B' गुण मिळालेला एकमेव विषय होता. मात्र त्यांच्या माध्यमिक शाळेच्या नोंदींनुसार तर्कशास्त्रासह सर्व विषयांत त्यांना 'A' गुण मिळाले होते. त्यांनी १९१८ ते १९२४ या काळात वॉर्सा विद्यापीठात शिक्षण घेतले. सुरुवातीला त्यांचा जीवशास्त्र शिकण्याचा विचार होता; पण विद्यापीठ हे वॉर्सा तर्कशास्त्र व तत्त्वज्ञान परंपरेचे केंद्र असल्याने त्यांना गणित, तत्त्वज्ञान आणि तर्कशास्त्रात रस निर्माण झाला. १९२४ मध्ये त्यांनी स्तानिस्वाव लेश्न्येव्स्की (Stanis\l{}aw Le\'{s}niewski) यांच्या मार्गदर्शनाखाली डॉक्टरेट मिळवली.

१९३९ मध्ये अमेरिकेत स्थलांतर करण्यापूर्वी टार्स्की यांनी वॉर्सा येथे माध्यमिक शाळेत शिकवत असताना आपले काही सर्वात महत्त्वाचे काम पूर्ण केले. तार्किक परिणाम आणि तार्किक सत्य यांवरील त्यांचे लेख याच काळात लिहिले गेले. १९३९ मध्ये ते व्याख्यानांच्या दौऱ्यासाठी अमेरिकेत गेले होते. त्याच भेटीदरम्यान जर्मनीने पोलंडवर आक्रमण केले; ज्यू वंशामुळे टार्स्की परत जाऊ शकले नाहीत. त्यांची पत्नी आणि मुले युद्ध संपेपर्यंत पोलंडमध्ये राहिली; नंतर त्यांनाही अमेरिकेत स्थलांतर करता आले. टार्स्की यांनी हार्वर्ड, सिटी कॉलेज ऑफ न्यू यॉर्क, प्रिन्स्टन येथील इन्स्टिट्यूट फॉर ॲडव्हान्स्ड स्टडी आणि अखेरीस कॅलिफोर्निया विद्यापीठ, बर्कली येथे अध्यापन केले. बर्कलीमध्ये त्यांनी तर्कशास्त्र आणि विज्ञानपद्धती यांचा बहुविषयक अभ्यास कार्यक्रम स्थापन केला. २६ ऑक्टोबर १९८३ रोजी वयाच्या ८२व्या वर्षी टार्स्की यांचे निधन झाले.

\paragraph{अधिक वाचन.}
टार्स्की यांच्या जीवनाविषयी अधिक माहितीसाठी \emph{Alfred Tarski: Life
  and Logic} हे चरित्र (Feferman 2004) पाहा. तार्किक परिणाम आणि सत्य यांवरील टार्स्की यांचे मूलभूत लेख इंग्रजीत (Corcoran 1983) येथे उपलब्ध आहेत. त्यांचे सर्व मूळ लेख चार खंडांच्या मालिकेत संकलित केले आहेत (Tarski 1981).














\section{ॲलन ट्यूरिंग}\label{his:bio:tur:sec}

ॲलन ट्यूरिंग यांचा जन्म २३ जून १९१२ रोजी लंडनमधील मायडा व्हेल येथे झाला. त्यांना सैद्धांतिक संगणकशास्त्राचे जनक मानले जाते. भौतिक विज्ञान आणि गणित यांविषयी त्यांना लहानपणापासूनच रस होता. पण ते शिकत असलेल्या शाळांत साहित्य आणि अभिजात विषयांवर भर असल्याने त्यांच्या आवडीला फार वाव मिळाला नाही. परिणामी, शाळेत त्यांची कामगिरी कमकुवत राहिली आणि अनेक शिक्षकांनी त्यांना तंबी दिली.



ट्यूरिंग यांनी केंब्रिजच्या किंग्ज कॉलेजमध्ये गणिताचे पदवीशिक्षण घेतले. गणनीय फलन ही संकल्पना नेमकी परिभाषित करण्यासाठी आणि निर्णय समस्या अनिर्णेय आहे हे सिद्ध करण्यासाठी त्यांनी १९३६ मध्ये पुढे ट्यूरिंग यंत्र म्हणून ओळखली गेलेली संकल्पना विकसित केली. ॲलोन्झो चर्च यांनी स्वतःच्या लॅम्डा कलनाच्या साहाय्याने हा निष्कर्ष आधीच सिद्ध केला होता. तरीही चर्च यांच्या निष्कर्षाचा उल्लेख करून ट्यूरिंग यांचा लेख प्रकाशित झाला. चर्च यांच्या निमंत्रणावरून ट्यूरिंग प्रिन्स्टनला गेले. तेथे ते १९३६ ते १९३८ या काळात राहिले आणि चर्च यांच्या मार्गदर्शनाखाली डॉक्टरेट मिळवली.

तर्कशास्त्रातील रस असूनही भौतिक विज्ञानातील ट्यूरिंग यांची जुनी आवड कायम होती. दुसऱ्या महायुद्धात ब्रिटनच्या ब्लेचली पार्क येथील संकेतभेदन विभागात काम करताना त्यांच्या व्यावहारिक कौशल्याचा उपयोग झाला. जर्मन नौदलाच्या संदेशांमध्ये वापरला जाणारा एनिग्मा संकेत भेदण्यात त्यांची मध्यवर्ती भूमिका होती. संख्याशास्त्र आणि संकेतलेखनातील त्यांचे कौशल्य, तसेच विद्युत् आणि यांत्रिक तंत्रांचा वापर, यांमुळे चमूला “बॉम्ब” (bombe) नावाचे संकेतभेदक यंत्र तयार करून हा संकेत भेदता आला. त्यांच्या कल्पनांनी जगातील पहिला प्रोग्राम करता येणारा इलेक्ट्रॉनिक संगणक, कोलॉसस, तयार करण्यासही मदत केली; ब्लेचली पार्क येथेच तो जर्मन लोरेन्झ संकेत भेदण्यासाठी वापरला गेला.

ट्यूरिंग समलैंगिक होते. तरीही १९४१ मध्ये त्यांनी ब्लेचली पार्कमधील सहकारी जोन क्लार्क यांना लग्नाची मागणी घातली. नंतर त्यांनी साखरपुडा मोडला आणि आपण समलैंगिक आहोत असे त्यांना सांगितले. ट्यूरिंग यांचे आयुष्यात अनेक पुरुष जोडीदार होते; त्या काळी ब्रिटनमध्ये पुरुषांमधील समलैंगिक संबंध कायद्याने गुन्हा होते. १९५२ मध्ये त्यावेळच्या जोडीदाराच्या एका मित्राने ट्यूरिंग यांच्या घरात चोरी केली. पोलिसांत तक्रार करताना ट्यूरिंग यांनी आपले समलैंगिक संबंध असल्याचे सांगितले; अशा संबंधांना लवकरच कायदेशीर मान्यता मिळेल, असा त्यांचा समज होता. मात्र तसे झाले नाही आणि त्यांच्यावर “गंभीर अश्लीलता” (gross indecency) या तत्कालीन गुन्ह्याचा आरोप ठेवण्यात आला. तुरुंगवासाऐवजी ट्यूरिंग यांनी कामेच्छा कमी करणारे संप्रेरक उपचार निवडले. ८ जून १९५४ रोजी सायनाइडचे अतिप्रमाण सेवन झाल्याने ते मृतावस्थेत आढळले; स्रोताच्या मते ती बहुधा आत्महत्या असावी. २०१३ मध्ये राणी एलिझाबेथ दुसरी यांनी त्यांना राजमाफी दिली.

\paragraph{अधिक वाचन.}
ॲलन ट्यूरिंग यांच्या सविस्तर चरित्रासाठी Hodges (2014) पाहा. त्यांच्या जीवनावर आणि कार्यावर आधारित \emph{Breaking the Code} हे नाटक १९९६ मध्ये डेरेक जेकोबी यांनी ट्यूरिंग यांची भूमिका केलेल्या दूरचित्रवाणी कार्यक्रमाच्या रूपात सादर झाले. अकादमी पुरस्कारासाठी नामांकन मिळालेला \emph{The Imitation Game} हा बेनेडिक्ट कंबरबॅच आणि कीरा नाइटली यांच्या भूमिका असलेला चित्रपटही ट्यूरिंग यांचे जीवन आणि ब्लेचली पार्कमधील काळ यांवर सैलपणे आधारलेला आहे (Tyldum 2014).

Radiolab (2012) येथे ट्यूरिंग यांच्या जीवन आणि कार्यावरील अनेक पॉडकास्ट आहेत. बीबीसी होरायझनचा \emph{The Strange Life and Death of
  Dr.~Turing} हा माहितीपट ऑनलाइन पाहता येतो (Sykes 1992). ट्यूरिंग यांच्या जन्मशताब्दीनिमित्त २०१२ मध्ये तयार केलेले कार्यरत लेगो ट्यूरिंग यंत्र दाखवणारा छोटा व्हिडिओ (Theelen 2012) येथे आहे.

ट्यूरिंग यंत्र आणि निर्णय समस्येवरील ट्यूरिंग यांचा मूळ लेख Turing (1937) हा आहे.














\section{एर्न्स्ट झर्मेलो}\label{his:bio:zer:sec}



एर्न्स्ट झर्मेलो यांचा जन्म २७ जुलै १८७१ रोजी जर्मनीतील बर्लिन येथे झाला. त्यांना पाच बहिणी होत्या; मात्र कुटुंबातील सदस्यांची तब्येत नाजूक असल्याने त्यांपैकी केवळ तीनच जणी प्रौढ वयापर्यंत जगल्या. झर्मेलो लहान असतानाच त्यांच्या आईवडिलांचेही निधन झाले आणि वयाच्या सतराव्या वर्षी ते व त्यांची भावंडे अनाथ झाली. झर्मेलो यांना कलांत, विशेषतः कवितेत, मोठा रस होता. ते तीक्ष्ण बुद्धीचे, चतुर आणि टीकात्मक स्वभावाचे म्हणून ओळखले जात. त्यांच्या सर्वाधिक प्रसिद्ध गणितीय कामगिरीत १९०४ मध्ये निवडीचे स्वयंसिद्धक मांडणे आणि १९०८ मध्ये संच उपपत्तीची स्वयंसिद्धकीय मांडणी करणे यांचा समावेश होतो.

विद्यापीठात झर्मेलो यांच्या आवडी विविध होत्या. त्यांनी भौतिकशास्त्र, गणित आणि तत्त्वज्ञान यांचे अभ्यासक्रम घेतले. हेरमान श्वार्त्स यांच्या मार्गदर्शनाखाली त्यांनी १८९४ मध्ये बर्लिन विद्यापीठात \emph{Investigations in
  the Calculus of Variations} हा प्रबंध पूर्ण केला. १८९७ मध्ये पुढील शिक्षणासाठी त्यांनी गॉटिंगेन विद्यापीठात जाण्याचे ठरवले; तेथे डेव्हिड हिल्बर्ट यांच्या गणिताच्या पायांवरील कार्याचा त्यांच्यावर मोठा प्रभाव पडला. १८९९ मध्ये ते प्राध्यापकपदासाठी पात्र झाले, पण वेतन मिळणारे प्राध्यापकपद अकरा वर्षांनीच मिळाले. त्यांच्या विचित्र वागण्यामुळे आणि “अस्वस्थ उतावळेपणा”मुळे हा विलंब झाला असावा, असा स्रोताचा अंदाज आहे.

अखेरीस १९१० मध्ये झर्मेलो यांना झुरिक विद्यापीठात वेतन मिळणारे प्राध्यापकपद मिळाले; मात्र क्षयरोगामुळे १९१६ मध्ये त्यांना निवृत्त व्हावे लागले. बरे झाल्यावर १९२६ मध्ये त्यांना फ्रायबुर्ग विद्यापीठात मानद प्राध्यापकपद देण्यात आले. या काळात त्यांनी गणिताच्या पायांवर काम केले. थोराल्फ स्कोलेम आणि कुर्ट गोडेल यांच्या कार्यावर ते नाराज झाले आणि त्यांनी आपल्या लेखांत त्यांच्या दृष्टिकोनांवर जाहीर टीका केली. १९३५ मध्ये हिटलरला अभिवादन करण्यास नकार दिल्यामुळे त्यांच्याविरुद्ध तक्रार झाल्यानंतर त्यांना मानद प्राध्यापकपदाचा राजीनामा देणे भाग पडले.

झर्मेलो यांच्या आयुष्याचा उत्तरार्ध एकाकीपणात गेला. १९३५ मध्ये राजीनामा द्यावा लागल्यानंतर त्यांनी गणिताचे काम सोडले. ते ग्रामीण भागात जाऊन साधेपणाने राहू लागले. १९४४ मध्ये त्यांनी विवाह केला. त्यांची दृष्टी मंदावत असल्याने ते पत्नीवर पूर्णपणे अवलंबून झाले. १९५१ पर्यंत त्यांची दृष्टी पूर्णपणे गेली होती. २१ मे १९५३ रोजी जर्मनीतील ग्युंटरस्टाल येथे त्यांचे निधन झाले.

\paragraph{अधिक वाचन.}
झर्मेलो यांच्या सविस्तर चरित्रासाठी Ebbinghaus (2015) पाहा. त्यांचे १९०४ आणि १९०८ मधील मूलभूत लेख मूळ जर्मन भाषेत (Zermelo 1904; Zermelo 1908) येथे उपलब्ध आहेत. भौतिकशास्त्रावरील लेखनासह त्यांचे संकलित लेख इंग्रजी भाषांतरात (Ebbinghaus 2010; Ebbinghaus 2013) येथे उपलब्ध आहेत.




